\documentclass[11pt]{article}

\usepackage[utf8]{inputenc}
\usepackage[T1]{fontenc}
\usepackage{lmodern}
\usepackage{geometry}
\usepackage{amsmath,amssymb,amsfonts,amsthm,mathtools}
\usepackage{bm}
\usepackage{enumitem}
\usepackage{microtype}
\usepackage{hyperref}
\usepackage{titlesec}
\usepackage{booktabs}
\usepackage{array}
\usepackage{setspace}
\usepackage{mathrsfs}
\usepackage{physics}

\geometry{margin=1in}
\hypersetup{colorlinks=true, linkcolor=black, urlcolor=blue, citecolor=black}
\setlist[enumerate]{topsep=0.4em,itemsep=0.35em}
\setlist[itemize]{topsep=0.3em,itemsep=0.25em}
\titleformat{\section}{\large\bfseries}{\thesection.}{0.5em}{}
\titleformat{\subsection}{\normalsize\bfseries}{\thesubsection.}{0.5em}{}
\setstretch{1.06}

\newcommand{\Aether}{\AE{}ther}
\newcommand{\Aflow}{\AE{}ther-flow}
\newcommand{\TheoryName}{\Aflow{} Interpretation of Relativity}
\newcommand{\TheoryProgram}{\Aether{} / \Aflow{} framework}

\newcommand{\CompletedMetric}{g^{\sharp}}
\newcommand{\MatterSpace}{\Psi}

\newcommand{\Car}{\mathrm{car}}
\newcommand{\Cur}{\mathrm{cur}}
\newcommand{\Hom}{\mathrm{hom}}

\newcommand{\ForSpace}[1]{\mathcal I_{#1}^{\mathrm{for}}}
\newcommand{\PrimShell}{\mathcal S_{\mathrm{prim}}}
\newcommand{\MetricOut}{\mathscr G_{\mathrm{out}}}
\newcommand{\MatterOut}{\mathscr T_{\mathrm{out}}}

\newcommand{\SubSectionPackage}{\mathfrak V^{\mathrm{subsec}}}
\newcommand{\MicroSectionPackage}{\mathfrak W^{\mathrm{microsec}}}
\newcommand{\MicroSectionSpace}[1]{\mathcal N_{#1}^{\mathrm{microsec}}}
\newcommand{\MicroSectionBody}[1]{\mathcal B_{#1}^{\mathrm{microsec}}}
\newcommand{\MicroSectionFunctional}[1]{\mathscr E_{#1}^{\mathrm{microsec}}}
\newcommand{\MicroSectionShell}[1]{\mathcal Z_{#1}^{\mathrm{microsec}}}
\newcommand{\SubPreSectionCore}{\mathfrak U^{\mathrm{sec}}}

\newcommand{\NanoSectionPackage}{\mathfrak X^{\mathrm{nanosec}}}
\newcommand{\NanoSectionSpace}[1]{\mathcal P_{#1}^{\mathrm{nanosec}}}
\newcommand{\NanoSectionBody}[1]{\mathcal B_{#1}^{\mathrm{nanosec}}}
\newcommand{\NanoSectionProj}[1]{\Pi_{#1}^{\mathrm{nanosec}}}
\newcommand{\NanoSectionFiber}[2]{\mathcal F_{#1}^{\mathrm{nanosec}}(#2)}
\newcommand{\NanoSectionMicroMin}[1]{\mathfrak n_{#1}^{\mathrm{nanosec}}}
\newcommand{\NanoSectionFunctional}[1]{\mathscr N_{#1}^{\mathrm{nanosec}}}
\newcommand{\NanoSectionShell}[1]{\mathcal Z_{#1}^{\mathrm{nanosec}}}

\newcommand{\PicoSectionPackage}{\mathfrak Y^{\mathrm{picosec}}}
\newcommand{\PicoSectionSpace}[1]{\mathcal Q_{#1}^{\mathrm{picosec}}}
\newcommand{\PicoSectionBody}[1]{\mathcal B_{#1}^{\mathrm{picosec}}}
\newcommand{\PicoSectionProj}[1]{\Pi_{#1}^{\mathrm{picosec}}}
\newcommand{\PicoSectionFiber}[2]{\mathcal F_{#1}^{\mathrm{picosec}}(#2)}
\newcommand{\PicoSectionNanoMin}[1]{\mathfrak p_{#1}^{\mathrm{picosec}}}
\newcommand{\PicoSectionFunctional}[1]{\mathscr P_{#1}^{\mathrm{picosec}}}
\newcommand{\PicoSectionShell}[1]{\mathcal Z_{#1}^{\mathrm{picosec}}}

\newcommand{\PicoImageCore}{\mathfrak Y^{\mathrm{pic,img}}}
\newcommand{\PicoImgBody}[1]{\mathcal B_{#1}^{\mathrm{pic,img}}}
\newcommand{\PicoImgProj}[1]{\Pi_{#1}^{\mathrm{pic,img}}}
\newcommand{\PicoImgFunctional}[1]{\mathscr P_{#1}^{\mathrm{pic,img}}}
\newcommand{\PicoImgShell}[1]{\mathcal Z_{#1}^{\mathrm{pic,img}}}

\newtheorem{definition}{Definition}
\newtheorem{proposition}{Proposition}
\newtheorem{remark}{Remark}
\newtheorem{corollary}{Corollary}
\newtheorem{theorem}{Theorem}

\title{\textbf{The \TheoryName{}}\\[0.4em]
\large Primitive-Reservoir Observer-Reduction\\
\large Nano-Section-Generating Substrate Pico-Section Dynamics\\
\large Collapse to an Observer-Relevant Pico-Section Minimizer-Image Core}
\author{}
\date{}

\begin{document}

\maketitle

\begin{abstract}
The current primitive-reservoir observer-reduction frontier already moved the live burden below the micro-section-generating substrate nano-section dynamics package itself: the active pico-section theorem derived the current nano-section package \(\NanoSectionPackage\) canonically from a deeper nano-section-generating substrate pico-section dynamics package \(\PicoSectionPackage\). Under the recorded routing safeguard, however, the next honest move at that level is not another same-output deeper-origin relay that leaves \(\PicoSectionPackage\) unchanged. The sharper benchmark-facing gain is to prove a genuinely smaller theorem-level collapse strictly inside that current pico package.

The present note takes that route. It isolates the \emph{observer-relevant pico-section minimizer-image core}: the minimizer-image compact body over the recorded nano-section body, the restricted projection back to that body, the restricted pico-section functional on that image, the exact-shell image of the recorded nano-section shell, and benchmark faithfulness. The full exact pico-section shell is then shown to be forced by that minimizer image, and ambient pico-section state-space directions outside the minimizer image do not add independent benchmark-facing freedom once the induced nano-section package is held fixed.

The benchmark boundary remains fixed throughout. The one operative metric is still exactly \(\CompletedMetric\), the ordinary matter route is unchanged, there is no second operative metric, the low-energy matter law is not enlarged, and no stronger promotion language is licensed. The positive result is therefore narrower than another deeper relay and sharper than the full pico-section package: the live burden now collapses to a smaller observer-relevant pico-section minimizer-image core. What remains open is to derive or justify that smaller core from still more primitive substrate structure, or else prove a still sharper collapse strictly inside it.
\end{abstract}

\section{Introduction}

The active derivational program of \TheoryName{} is again constrained by two routing facts. First, the current active-primary theorem already removed the current nano-section package as the primitive stopping point by deriving it from a more primitive pico-section package. Second, the repository has already fixed a safeguard against recursive depth drift: once a benchmark-relevant datum is in hand, another deeper-origin relay that preserves that same datum is no longer honest front-line progress merely by being deeper.

That safeguard matters here. The accompanying pico-frontier routing note records that the next honest move below the current pico-section frontier is either to derive or justify the current pico-section package from still more primitive substrate structure, or else to prove a genuinely smaller theorem-level collapse strictly inside it. The present note takes that second route.

Inside the recorded pico-section package, the ambient state-space presentation and the off-image body directions are not an independent observer-relevant burden. Once attention is restricted to the minimizer image over the recorded nano-section body and to the exact-shell image of the recorded nano-section shell, the omitted ambient directions are forced or benchmark-neutral.

\paragraph{Framework and claim boundary.}
\TheoryProgram{} denotes the broader \Aether{} / \Aflow{} framework built on the claims that the \Aether{} is the underlying four-dimensional substrate of reality, the \Aflow{} is its intrinsic ordered motion, observed three-dimensional space is the local experiential slice of that deeper substrate, S-time is the experienced order of change arising from matter, light, and the \Aflow{}, and observed expansion is the three-dimensional appearance of deeper four-dimensional ordered motion. Within that framework, \TheoryName{} denotes the exact-closure theory in which the effective gravitational dynamics are adopted to be exactly Einsteinian with universal matter coupling. In the active sequence, \emph{adoption} means use of that established relativistic dynamics without claiming substrate derivation, while \emph{derivation} is reserved for a first-principles recovery from explicit substrate variables.

The present note stays strictly inside that boundary. It does not introduce a second operative metric, it does not enlarge the low-energy matter law, and it does not reopen older screened layers as if they were again the live burden. Its positive role is narrower and benchmark facing: compress the live pico-section burden to the part of the package that still matters observer-relevantly.

\section{Fixed Inputs}

\begin{definition}[Recorded primitive continuation data]
\label{def:fixed}
Fix the following continuation data from the active primitive-reservoir observer-reduction line.
\begin{enumerate}
    \item For each sector label \(\sigma\in\{\Car,\Cur,\Hom\}\), a primitive forcing shell
    \[
    \ForSpace{\sigma}.
    \]
    \item The product shell
    \[
    \PrimShell
    \equiv
    \ForSpace{\Car}\times \ForSpace{\Cur}\times \ForSpace{\Hom}.
    \]
    \item The recorded metric and ordinary-matter outputs
    \[
    \MetricOut:\PrimShell\to\{\text{observer metrics}\},
    \qquad
    \MatterOut:\PrimShell\times\MatterSpace\to\{\text{observer stress tensors}\},
    \]
    satisfying
    \begin{equation}
    \MetricOut(\mathbf w)=\CompletedMetric,
    \qquad
    \MatterOut(\mathbf w,\psi)
    =
    T_{\mu\nu}^{\mathrm{matter}}[\CompletedMetric,\psi]
    \label{eq:fixedoutputs}
    \end{equation}
    for every \(\mathbf w\in\PrimShell\) and every matter configuration \(\psi\in\MatterSpace\).
\end{enumerate}
\end{definition}

\begin{definition}[Recorded nano-section downstream chain]
\label{def:downstream}
Fix \(\sigma\in\{\Car,\Cur,\Hom\}\). The active line already fixes the following downstream data and consequences:
\begin{enumerate}
    \item the recorded micro-section target data
    \[
    \MicroSectionSpace{\sigma},\qquad
    \MicroSectionBody{\sigma},\qquad
    \MicroSectionFunctional{\sigma},\qquad
    \MicroSectionShell{\sigma};
    \]
    \item the current exact-shell-transport-section-generating substrate sub-section dynamics package \(\SubSectionPackage\), the current micro-section package \(\MicroSectionPackage\), and the observer-relevant sub-pre-core exact-shell transport-section core \(\SubPreSectionCore\);
    \item the previously recorded fact that, once the recorded nano-section package is available, the current micro-section package, the current sub-section package, the observer-relevant section core, and the downstream shell-transport consequences used on the active line follow unchanged.
\end{enumerate}
\end{definition}

\begin{definition}[Recorded micro-section-generating substrate nano-section dynamics]
\label{def:nanosec}
Fix \(\sigma\in\{\Car,\Cur,\Hom\}\). The recorded current nano-section dynamics package \(\NanoSectionPackage\) for sector \(\sigma\) consists of the following data.
\begin{enumerate}
    \item A finite-dimensional Euclidean nano-section state space \(\NanoSectionSpace{\sigma}\), a compact convex nano-section body
    \[
    \NanoSectionBody{\sigma}\subset\NanoSectionSpace{\sigma}
    \]
    with nonempty interior, and an affine surjection
    \[
    \NanoSectionProj{\sigma}:\NanoSectionSpace{\sigma}\to\MicroSectionSpace{\sigma}
    \]
    satisfying
    \[
    \NanoSectionProj{\sigma}\bigl(\NanoSectionBody{\sigma}\bigr)=\MicroSectionBody{\sigma}.
    \]
    For each \(u\in\MicroSectionBody{\sigma}\), write
    \[
    \NanoSectionFiber{\sigma}{u}
    \equiv
    \bigl(\NanoSectionProj{\sigma}\bigr)^{-1}(u)\cap\NanoSectionBody{\sigma}.
    \]
    Each such fiber is required to be nonempty, compact, and convex.
    \item A nonnegative smooth nano-section functional
    \[
    \NanoSectionFunctional{\sigma}:\NanoSectionSpace{\sigma}\to[0,\infty)
    \]
    such that, for every \(u\in\MicroSectionBody{\sigma}\), the restriction of \(\NanoSectionFunctional{\sigma}\) to \(\NanoSectionFiber{\sigma}{u}\) is strictly convex and attains a unique minimizer. The resulting minimizer depends smoothly on \(u\); write it as
    \[
    \NanoSectionMicroMin{\sigma}:\MicroSectionBody{\sigma}\to\NanoSectionBody{\sigma}.
    \]
    These data satisfy
    \[
    \NanoSectionProj{\sigma}\circ\NanoSectionMicroMin{\sigma}
    =
    \mathrm{id}_{\MicroSectionBody{\sigma}},
    \]
    and the effective minimum-value map recovers the recorded micro-section functional:
    \[
    \MicroSectionFunctional{\sigma}(u)
    =
    \NanoSectionFunctional{\sigma}\bigl(\NanoSectionMicroMin{\sigma}(u)\bigr)
    =
    \min_{p\in\NanoSectionFiber{\sigma}{u}}
    \NanoSectionFunctional{\sigma}(p)
    \qquad
    \text{for every }
    u\in\MicroSectionBody{\sigma}.
    \]
    \item A closed embedded exact nano-section shell
    \[
    \NanoSectionShell{\sigma}\subset\NanoSectionBody{\sigma}
    \]
    satisfying
    \[
    \NanoSectionShell{\sigma}
    =
    \left\{
    p\in\NanoSectionBody{\sigma}:
    \begin{array}{l}
    \NanoSectionProj{\sigma}(p)\in\MicroSectionShell{\sigma},\\[0.2em]
    \NanoSectionFunctional{\sigma}(p)
    =
    \displaystyle\min_{q\in\NanoSectionFiber{\sigma}{\NanoSectionProj{\sigma}(p)}}
    \NanoSectionFunctional{\sigma}(q)
    \end{array}
    \right\},
    \]
    and
    \[
    \NanoSectionProj{\sigma}\big|_{\NanoSectionShell{\sigma}}
    :
    \NanoSectionShell{\sigma}
    \xrightarrow{\cong}
    \MicroSectionShell{\sigma}
    \]
    is a diffeomorphism.
    \item Benchmark faithfulness, meaning preservation of the benchmark outputs \eqref{eq:fixedoutputs}, no second operative metric, no enlarged low-energy matter law, and presentation as the fixed benchmark-relevant input for the current downstream chain.
\end{enumerate}
\end{definition}

\begin{definition}[Recorded pico-section dynamics]
\label{def:picosec}
Fix \(\sigma\in\{\Car,\Cur,\Hom\}\). The recorded current pico-section dynamics package \(\PicoSectionPackage\) for sector \(\sigma\) consists of the following data.
\begin{enumerate}
    \item A finite-dimensional Euclidean pico-section state space \(\PicoSectionSpace{\sigma}\), a compact convex pico-section body
    \[
    \PicoSectionBody{\sigma}\subset\PicoSectionSpace{\sigma}
    \]
    with nonempty interior, and an affine surjection
    \[
    \PicoSectionProj{\sigma}:\PicoSectionSpace{\sigma}\to\NanoSectionSpace{\sigma}
    \]
    satisfying
    \[
    \PicoSectionProj{\sigma}\bigl(\PicoSectionBody{\sigma}\bigr)=\NanoSectionBody{\sigma}.
    \]
    For each \(p\in\NanoSectionBody{\sigma}\), write
    \[
    \PicoSectionFiber{\sigma}{p}
    \equiv
    \bigl(\PicoSectionProj{\sigma}\bigr)^{-1}(p)\cap\PicoSectionBody{\sigma}.
    \]
    Each such fiber is required to be nonempty, compact, and convex.
    \item A nonnegative smooth pico-section functional
    \[
    \PicoSectionFunctional{\sigma}:\PicoSectionSpace{\sigma}\to[0,\infty)
    \]
    such that, for every \(p\in\NanoSectionBody{\sigma}\), the restriction of \(\PicoSectionFunctional{\sigma}\) to \(\PicoSectionFiber{\sigma}{p}\) is strictly convex and attains a unique minimizer. The resulting minimizer depends smoothly on \(p\); write it as
    \[
    \PicoSectionNanoMin{\sigma}:\NanoSectionBody{\sigma}\to\PicoSectionBody{\sigma}.
    \]
    These data satisfy
    \[
    \PicoSectionProj{\sigma}\circ\PicoSectionNanoMin{\sigma}
    =
    \mathrm{id}_{\NanoSectionBody{\sigma}},
    \]
    and the effective minimum value recovers the recorded nano-section functional:
    \[
    \NanoSectionFunctional{\sigma}(p)
    =
    \PicoSectionFunctional{\sigma}\bigl(\PicoSectionNanoMin{\sigma}(p)\bigr)
    =
    \min_{r\in\PicoSectionFiber{\sigma}{p}}
    \PicoSectionFunctional{\sigma}(r)
    \qquad
    \text{for every }
    p\in\NanoSectionBody{\sigma}.
    \]
    \item A closed embedded exact pico-section shell
    \[
    \PicoSectionShell{\sigma}\subset\PicoSectionBody{\sigma}
    \]
    satisfying
    \[
    \PicoSectionShell{\sigma}
    =
    \left\{
    r\in\PicoSectionBody{\sigma}:
    \begin{array}{l}
    \PicoSectionProj{\sigma}(r)\in\NanoSectionShell{\sigma},\\[0.2em]
    \PicoSectionFunctional{\sigma}(r)
    =
    \displaystyle\min_{s\in\PicoSectionFiber{\sigma}{\PicoSectionProj{\sigma}(r)}}
    \PicoSectionFunctional{\sigma}(s)
    \end{array}
    \right\},
    \]
    and
    \[
    \PicoSectionProj{\sigma}\big|_{\PicoSectionShell{\sigma}}
    :
    \PicoSectionShell{\sigma}
    \xrightarrow{\cong}
    \NanoSectionShell{\sigma}
    \]
    is a diffeomorphism.
    \item Benchmark faithfulness, meaning preservation of the benchmark outputs \eqref{eq:fixedoutputs}, no second operative metric, no enlarged low-energy matter law, and presentation as a deeper substrate-side source for the current nano-section package rather than as a replacement benchmark theory.
\end{enumerate}
\end{definition}

\begin{remark}
Definitions \ref{def:nanosec} and \ref{def:picosec} fix the current frontier package. The present note does not spend its move on another deeper derivation below that frontier. It asks instead whether the full recorded pico-section package carries benchmark-neutral ambient freedom that can be removed without changing the induced nano-section package or the downstream chain that depends on it.
\end{remark}

\section{Observer-Relevant Pico-Section Minimizer-Image Core}

\begin{definition}[Observer-relevant pico-section minimizer-image core]
\label{def:imagecore}
Fix \(\sigma\in\{\Car,\Cur,\Hom\}\) and a benchmark-faithful pico-section package in the sense of Definition \ref{def:picosec}. The \emph{observer-relevant pico-section minimizer-image core} \(\PicoImageCore\) for sector \(\sigma\) consists of:
\begin{enumerate}
    \item the compact minimizer-image body
    \[
    \PicoImgBody{\sigma}
    \equiv
    \PicoSectionNanoMin{\sigma}\bigl(\NanoSectionBody{\sigma}\bigr)
    \subset
    \PicoSectionBody{\sigma},
    \]
    together with the restricted projection
    \[
    \PicoImgProj{\sigma}
    \equiv
    \PicoSectionProj{\sigma}\big|_{\PicoImgBody{\sigma}}
    :
    \PicoImgBody{\sigma}\to\NanoSectionBody{\sigma};
    \]
    \item the restricted functional
    \[
    \PicoImgFunctional{\sigma}
    \equiv
    \PicoSectionFunctional{\sigma}\big|_{\PicoImgBody{\sigma}};
    \]
    \item the exact-shell image
    \[
    \PicoImgShell{\sigma}
    \equiv
    \PicoSectionNanoMin{\sigma}\bigl(\NanoSectionShell{\sigma}\bigr)
    \subset
    \PicoSectionShell{\sigma};
    \]
    \item benchmark faithfulness in the sense of Definition \ref{def:picosec}(4).
\end{enumerate}
\end{definition}

\begin{remark}
Definition \ref{def:imagecore} is intentionally smaller than Definition \ref{def:picosec}. It keeps only the minimizer image over the recorded nano-section body together with the restricted functional and the exact-shell image it already determines. The ambient pico-section state space and off-image body directions are not taken as independent core data.
\end{remark}

\section{Collapse of the Full Pico-Section Package}

\begin{proposition}[The minimizer-image body carries the body-level content]
\label{prop:body}
For every benchmark-faithful pico-section package, \(\PicoImgBody{\sigma}\) is compact and the restricted projection
\[
\PicoImgProj{\sigma}:\PicoImgBody{\sigma}\to\NanoSectionBody{\sigma}
\]
is a homeomorphism with inverse \(\PicoSectionNanoMin{\sigma}\).
\end{proposition}

\begin{proof}
The body \(\NanoSectionBody{\sigma}\) is compact by Definition \ref{def:nanosec}, and \(\PicoSectionNanoMin{\sigma}\) is continuous by Definition \ref{def:picosec}(2). Therefore
\[
\PicoImgBody{\sigma}
=
\PicoSectionNanoMin{\sigma}\bigl(\NanoSectionBody{\sigma}\bigr)
\]
is compact. The identity
\[
\PicoSectionProj{\sigma}\circ\PicoSectionNanoMin{\sigma}
=
\mathrm{id}_{\NanoSectionBody{\sigma}}
\]
shows that \(\PicoImgProj{\sigma}\) is surjective and that \(\PicoSectionNanoMin{\sigma}\) is a right inverse for it. If \(r_{1},r_{2}\in\PicoImgBody{\sigma}\) satisfy
\[
\PicoImgProj{\sigma}(r_{1})=\PicoImgProj{\sigma}(r_{2})=p,
\]
then each \(r_{i}\) lies in the fiber \(\PicoSectionFiber{\sigma}{p}\) and is the unique minimizer there, so \(r_{1}=r_{2}=\PicoSectionNanoMin{\sigma}(p)\). Thus \(\PicoImgProj{\sigma}\) is bijective with inverse \(\PicoSectionNanoMin{\sigma}\). Since \(\PicoImgBody{\sigma}\) is compact and \(\NanoSectionBody{\sigma}\) is Hausdorff, the continuous bijection \(\PicoImgProj{\sigma}\) is a homeomorphism.
\end{proof}

\begin{proposition}[The recorded nano-section functional is carried by the image core]
\label{prop:functional}
For every benchmark-faithful pico-section package, the recorded nano-section functional is recovered entirely from the restricted image-core functional:
\[
\NanoSectionFunctional{\sigma}(p)
=
\PicoImgFunctional{\sigma}\bigl(\PicoSectionNanoMin{\sigma}(p)\bigr)
\qquad
\text{for every }p\in\NanoSectionBody{\sigma}.
\]
\end{proposition}

\begin{proof}
By Definition \ref{def:imagecore}(2),
\[
\PicoImgFunctional{\sigma}\bigl(\PicoSectionNanoMin{\sigma}(p)\bigr)
=
\PicoSectionFunctional{\sigma}\bigl(\PicoSectionNanoMin{\sigma}(p)\bigr).
\]
Definition \ref{def:picosec}(2) then gives
\[
\PicoSectionFunctional{\sigma}\bigl(\PicoSectionNanoMin{\sigma}(p)\bigr)
=
\NanoSectionFunctional{\sigma}(p)
\]
for every \(p\in\NanoSectionBody{\sigma}\).
\end{proof}

\begin{proposition}[The full exact pico-section shell is forced by the image core]
\label{prop:shell}
For every benchmark-faithful pico-section package,
\[
\PicoSectionShell{\sigma}
=
\PicoImgShell{\sigma}
=
\PicoSectionNanoMin{\sigma}\bigl(\NanoSectionShell{\sigma}\bigr),
\]
and the restricted projection
\[
\PicoImgProj{\sigma}\big|_{\PicoImgShell{\sigma}}
:
\PicoImgShell{\sigma}\xrightarrow{\cong}\NanoSectionShell{\sigma}
\]
is a diffeomorphism.
\end{proposition}

\begin{proof}
If \(r\in\PicoSectionShell{\sigma}\), then Definition \ref{def:picosec}(3) gives
\[
p\equiv \PicoSectionProj{\sigma}(r)\in\NanoSectionShell{\sigma},
\]
and \(r\) minimizes \(\PicoSectionFunctional{\sigma}\) on the fiber \(\PicoSectionFiber{\sigma}{p}\). By Definition \ref{def:picosec}(2), that minimizer is unique and equals \(\PicoSectionNanoMin{\sigma}(p)\). Hence
\[
r=\PicoSectionNanoMin{\sigma}(p)\in\PicoImgShell{\sigma},
\]
so \(\PicoSectionShell{\sigma}\subset\PicoImgShell{\sigma}\).

Conversely, if \(p\in\NanoSectionShell{\sigma}\), then \(\PicoSectionNanoMin{\sigma}(p)\) is by definition the unique minimizer of \(\PicoSectionFunctional{\sigma}\) on the fiber over \(p\). Therefore \(\PicoSectionNanoMin{\sigma}(p)\in\PicoSectionShell{\sigma}\), so \(\PicoImgShell{\sigma}\subset\PicoSectionShell{\sigma}\). This proves the equality of shells.

Finally,
\[
\PicoImgProj{\sigma}\big|_{\PicoImgShell{\sigma}}
=
\PicoSectionProj{\sigma}\big|_{\PicoSectionShell{\sigma}},
\]
and the latter is a diffeomorphism onto \(\NanoSectionShell{\sigma}\) by Definition \ref{def:picosec}(3).
\end{proof}

\begin{proposition}[The image core canonically reproduces the current nano-section package]
\label{prop:package}
Fix a benchmark-faithful pico-section package. Then the current nano-section package of Definition \ref{def:nanosec} is exactly the benchmark-relevant package induced by the smaller image core. In particular:
\begin{enumerate}
    \item the recorded nano-section body and projection are recovered by Proposition \ref{prop:body};
    \item the recorded nano-section functional on \(\NanoSectionBody{\sigma}\) is recovered by Proposition \ref{prop:functional};
    \item the recorded exact nano-section shell is recovered by Proposition \ref{prop:shell}.
\end{enumerate}
\end{proposition}

\begin{proof}
Item (1) is Proposition \ref{prop:body}. Item (2) is Proposition \ref{prop:functional}. Item (3) is Proposition \ref{prop:shell}. These statements are exactly the benchmark-relevant constituents of the current nano-section package.
\end{proof}

\begin{proposition}[All downstream recovery factors through the image core]
\label{prop:downstream}
Fix a benchmark-faithful pico-section package. Then all downstream recovery of the current micro-section package \(\MicroSectionPackage\), the current sub-section package \(\SubSectionPackage\), the observer-relevant sub-pre-core exact-shell transport-section core \(\SubPreSectionCore\), and the previously recorded shell-transport consequences on the active line factors through \(\PicoImageCore\).
\end{proposition}

\begin{proof}
Definition \ref{def:downstream}(3) records that the active downstream chain depends on the recorded nano-section package. Proposition \ref{prop:package} shows that this entire nano-section package is already determined by the smaller image core. Therefore the downstream chain does not require the full ambient pico-section package as additional independent data; it factors through \(\PicoImageCore\).
\end{proof}

\begin{proposition}[Ambient pico-section directions are benchmark neutral at current scope]
\label{prop:neutral}
For every benchmark-faithful pico-section package, the ambient pico-section state space \(\PicoSectionSpace{\sigma}\) and the off-image body directions
\[
\PicoSectionBody{\sigma}\setminus\PicoImgBody{\sigma}
\]
do not add independent benchmark-facing freedom beyond \(\PicoImageCore\).
\end{proposition}

\begin{proof}
By Proposition \ref{prop:body}, the full body-level content relevant to the recorded nano-section package is already fixed by the minimizer-image body \(\PicoImgBody{\sigma}\) and the restricted projection \(\PicoImgProj{\sigma}\). By Proposition \ref{prop:functional}, the current nano-section functional is already fixed by the restricted image-core functional \(\PicoImgFunctional{\sigma}\). By Proposition \ref{prop:shell}, the full exact pico-section shell is not additional data beyond the shell image; it is exactly that image. By Proposition \ref{prop:downstream}, all downstream recovery of the current micro-section package, the current sub-section package, the observer-relevant section core, and the previously recorded shell-transport consequences factors through this smaller core. Finally, benchmark faithfulness is already part of Definition \ref{def:imagecore}(4) and preserves the fixed outputs \eqref{eq:fixedoutputs}. Therefore nothing benchmark-facing at the present frontier requires the ambient pico-section state-space presentation or body points outside the minimizer image as additional independent data.
\end{proof}

\begin{proposition}[The benchmark package remains fixed]
\label{prop:benchmark}
Every observer-relevant pico-section minimizer-image core preserves the same benchmark one-metric package \eqref{eq:fixedoutputs}. In particular, the operative metric remains exactly \(\CompletedMetric\), the ordinary matter route is unchanged, there is no second operative metric, and the low-energy matter law is not enlarged.
\end{proposition}

\begin{proof}
This is part of benchmark faithfulness in Definition \ref{def:imagecore}(4). The present note only removes benchmark-neutral ambient pico-section freedom from the live burden. It does not alter the benchmark exact-closure package.
\end{proof}

\section{Main Theorem}

\begin{theorem}[Nano-section-generating substrate pico-section dynamics collapse to an observer-relevant pico-section minimizer-image core]
\label{thm:main}
Fix the primitive continuation data of Definition \ref{def:fixed}, the recorded nano-section downstream chain of Definition \ref{def:downstream}, the recorded nano-section package of Definition \ref{def:nanosec}, and a benchmark-faithful pico-section package in the sense of Definition \ref{def:picosec}. Then, for each sector \(\sigma\in\{\Car,\Cur,\Hom\}\), the live benchmark-facing burden inside the current pico-section frontier collapses from the full pico-section package \(\PicoSectionPackage\) to the smaller observer-relevant pico-section minimizer-image core \(\PicoImageCore\). More precisely:
\begin{enumerate}
    \item the body-level content of the full pico-section package is already carried by the minimizer-image body \(\PicoImgBody{\sigma}\), the restricted projection \(\PicoImgProj{\sigma}\), and the restricted image-core functional \(\PicoImgFunctional{\sigma}\) by Propositions \ref{prop:body} and \ref{prop:functional};
    \item the full exact pico-section shell is forced to equal the minimizer image of the recorded nano-section shell by Proposition \ref{prop:shell};
    \item the current nano-section package and therefore the current micro-section package \(\MicroSectionPackage\), the current sub-section package \(\SubSectionPackage\), the observer-relevant section core \(\SubPreSectionCore\), and the previously recorded downstream consequences all factor through \(\PicoImageCore\) by Propositions \ref{prop:package} and \ref{prop:downstream};
    \item ambient pico-section directions outside that image core are benchmark neutral at current scope by Proposition \ref{prop:neutral};
    \item the benchmark boundary remains fixed by Proposition \ref{prop:benchmark};
    \item consequently the remaining honest burden below the previous pico-section frontier is no longer the full nano-section-generating substrate pico-section dynamics package \(\PicoSectionPackage\) as such. What remains is to derive or justify the sharper image core \(\PicoImageCore\) from still more primitive substrate structure, or else prove a still smaller theorem-level collapse strictly inside that image core.
\end{enumerate}
\end{theorem}

\begin{proof}
Item (1) is Propositions \ref{prop:body} and \ref{prop:functional}. Item (2) is Proposition \ref{prop:shell}. Item (3) is Propositions \ref{prop:package} and \ref{prop:downstream}. Item (4) is Proposition \ref{prop:neutral}. Item (5) is Proposition \ref{prop:benchmark}. Item (6) is the routing consequence of the previous items together with the active safeguard against counting another same-output deeper-origin relay as new front-line progress when the live benchmark-relevant datum is unchanged.
\end{proof}

\begin{corollary}[What must be done next]
\label{cor:next}
After Theorem \ref{thm:main}, the next honest benchmark-facing manuscript must either:
\begin{enumerate}
    \item derive or justify the observer-relevant pico-section minimizer-image core \(\PicoImageCore\) from still more primitive substrate structure in a way that changes benchmark-facing status;
    \item identify a genuinely smaller theorem-level observer-relevant datum strictly inside \(\PicoImageCore\) and prove a sharper collapse to it;
    \item do either of the previous two while preserving the same benchmark one-metric package and the same claim boundary.
\end{enumerate}
Another same-output deeper-origin relay that preserves this image core, another re-expansion back to the full pico-section package as if its screened ambient directions were still live, another reopening of the already-screened transport-image, sub-pre-core, or pre-core packages as if they were again the active burden, or any move that introduces a second operative metric, enlarges the ordinary matter law, or uses stronger promotion language does not satisfy the present burden. If a quantum continuation is pursued, it matters benchmark-facingly only insofar as it derives this same image core, collapses it further, or closes a benchmark derivation gate in a genuinely new way.
\end{corollary}

\begin{proof}
Theorem \ref{thm:main} shows that the full pico-section package is no longer the minimal benchmark-facing datum at this stage. Therefore a subsequent manuscript must improve on the smaller image core rather than merely restating the larger package.
\end{proof}

\section{Conclusion}

The positive result of this note is a disciplined compression step inside the current pico-section frontier. The previous active-primary theorem already removed the current nano-section package as the primitive stopping point by deriving it from a deeper pico-section package. The present theorem then takes the remaining honest internal route available there: it shows that the full pico-section package is itself not the minimal benchmark-facing datum.

What survives as the live burden is the observer-relevant pico-section minimizer-image core. That sharper datum consists of the minimizer-image compact body over the recorded nano-section body, the restricted projection back to that body, the restricted pico-section functional on that image, the exact-shell image of the recorded nano-section shell, and benchmark faithfulness. The full exact pico-section shell is forced to equal that image shell, and ambient pico-section directions outside the minimizer image do not add independent benchmark-facing freedom once the induced nano-section package is fixed.

This is the honest next burden after the present theorem. A new primary manuscript must derive or justify that sharper image core from still more primitive substrate structure, or else collapse it further, while preserving the same benchmark one-metric package and the same claim boundary. The current result does not yet complete a first-principles derivation of GR from substrate microphysics, but it does narrow the live derivational burden one level further on the active line.

\end{document}
